% =============================================================================
% File:    chord-substitution-en.tex
% Title:   The chord substitution closes the continuous Jensen equation
% Subtitle: A teaching note on three vestigial regularity hypotheses
% Author:  [AUTHOR]
% Venue:   Comptes Rendus Mathematique (Academie des sciences)
% Source:  satellites/o3-chord-substitution/09-draft0.3-manuscript.md
% Date:    2026-06-06
%
% Compilation:
%   pdflatex chord-substitution-en
%   bibtex   chord-substitution-en
%   pdflatex chord-substitution-en
%   pdflatex chord-substitution-en
%
% Class:   The Centre Mersenne `crmath` class is the official template
%          (download pack_author-crmath.zip from centre-mersenne.org/pack_author/).
%          A fallback amsart configuration is provided below for users without
%          the crmath class installed: comment out \documentclass{crmath} and
%          uncomment the amsart block.
% =============================================================================

% --- OFFICIAL CRAS CLASS (preferred) ---
% \documentclass{crmath}

% --- FALLBACK AMSART CONFIGURATION (active by default) ---
\documentclass[11pt,a4paper]{amsart}
\usepackage[utf8]{inputenc}
\usepackage[T1]{fontenc}
\usepackage{amssymb,amsmath,amsthm}
\usepackage{microtype}
\usepackage{booktabs}
\usepackage{array}
\usepackage{tabularx}
\usepackage[hidelinks]{hyperref}
\usepackage[english]{babel}
\usepackage{lineno}
\linenumbers
\usepackage[a4paper,margin=2.5cm]{geometry}
\renewcommand{\baselinestretch}{1.5} % double-spaced for submission (CRAS requirement)
\setlength{\parindent}{1em}
\setlength{\parskip}{0.25em}

% Roman top-level section numbering (CRAS Mathematique convention)
\renewcommand{\thesection}{\Roman{section}}
\renewcommand{\thesubsection}{\Roman{section}.\arabic{subsection}}

% Theorem environments
\newtheorem{theorem}{Theorem}
\newtheorem{corollary}[theorem]{Corollary}
\newtheorem{proposition}[theorem]{Proposition}
\newtheorem{lemma}[theorem]{Lemma}
\newtheorem{remark}[theorem]{Remark}

% Custom commands
\newcommand{\R}{\mathbb R}
\newcommand{\Q}{\mathbb Q}
\newcommand{\N}{\mathbb N}
\newcommand{\E}{\mathbb E}
\newcommand{\PR}{\mathbb P}
% --- END FALLBACK ---

\begin{document}

\title[The chord substitution closes the continuous Jensen equation]%
{The chord substitution closes the continuous Jensen equation\\
\normalsize\textit{A teaching note on three vestigial regularity hypotheses}}

\author{[AUTHOR]\thanks{Corresponding author. Affiliation:
[INSTITUTION, COUNTRY]. E-mail: [author@institution.country].}}

\date{\today}

\subjclass[2020]{39B22, 39B05}
\keywords{Jensen equation, Cauchy equation, Hamel basis, functional
equation, affine function, chord substitution}

\begin{abstract}
The continuous-coefficient form of the Jensen functional equation on a
real interval --- the form in which the coefficient varies over the full
continuum $[0,1]$ rather than over a single value such as $\tfrac12$
--- shares no Hamel-basis pathology with its discrete-coefficient sibling,
and is solved in closed form by a one-line chord substitution. Concretely:
for $G\colon [0,M]\to\R$ satisfying $pG(u_1)+(1-p)G(u_2)=G(pu_1+(1-p)u_2)$
for all $u_1,u_2\in[0,M]$ and all $p\in[0,1]$, $G$ is affine on $[0,M]$
--- meaning $G(v)=av+b$ for some $a,b\in\R$ --- with no measurability,
boundedness, or continuity hypothesis on~$G$. We give the proof, exhibit
the explicit affine formula, articulate the dictionary of three classically
required regularity hypotheses that become vestigial under the
continuous-coefficient form, and document one recurrence of the trap in the
surrogate-calibration literature together with the structural reason the
trap is predictable to recur in any derivation that saturates Jensen's
inequality across a wide class of two-point distributions.
\end{abstract}

% Bilingual metadata: French translation of title, abstract, keywords for the
% mandatory bilingual metadata block. Detailed French manuscript is in
% chord-substitution-fr.tex (separate file).
\renewcommand{\abstractname}{R\'esum\'e}
\begin{abstract}
\selectlanguage{french}
La forme \`a coefficient continu de l'\'equation fonctionnelle de Jensen
sur un intervalle r\'eel --- celle o\`u le coefficient parcourt tout le
continuum $[0,1]$ plut\^ot qu'une valeur unique telle que $\tfrac12$ ---
ne partage aucune pathologie de type base de Hamel avec son analogue \`a
coefficient discret, et se r\'esout en forme close par une substitution par
corde tenant en une ligne. Plus pr\'ecis\'ement~: pour
$G\colon[0,M]\to\R$ satisfaisant
$p\,G(u_1)+(1-p)\,G(u_2)=G(pu_1+(1-p)u_2)$ pour tous $u_1,u_2\in[0,M]$ et
tout $p\in[0,1]$, la fonction $G$ est affine sur $[0,M]$
--- c.-\`a-d. de la forme $G(v)=av+b$ avec $a,b\in\R$ --- sans qu'aucune
hypoth\`ese de mesurabilit\'e, de bornitude ou de continuit\'e ne soit
requise. Nous en donnons la d\'emonstration, exhibons la formule affine
explicite, articulons le dictionnaire des trois hypoth\`eses de
r\'egularit\'e classiquement requises qui deviennent vestigiales sous la
forme \`a coefficient continu, et documentons une occurrence du pi\`ege
dans la litt\'erature sur la calibration par perte de substitution.
\selectlanguage{english}

\textbf{Mots-cl\'es :} \'equation de Jensen~; \'equation de Cauchy~; base de
Hamel~; \'equation fonctionnelle~; fonction affine~; substitution par
corde.
\end{abstract}

\maketitle

% =============================================================================
\section{Introduction}
\label{sec:introduction}

The equation
\begin{equation}
p\,G(u_1) + (1-p)\,G(u_2) \;=\; G\bigl(p\,u_1 + (1-p)\,u_2\bigr)
\qquad (u_1,u_2\in I,\ p\in[0,1])
\tag{$\star$}\label{eq:star}
\end{equation}
on a real interval $I$ --- where $G\colon I\to\R$ is the unknown and
$u_1,u_2,p$ are the arguments and coefficient --- has at least three faces
in classical analysis, and they have been steadily conflated for over a
century.

The \textbf{discrete-coefficient} Jensen equation,
\begin{equation}
G\!\left(\tfrac{u_1+u_2}{2}\right) = \tfrac{G(u_1)+G(u_2)}{2},
\tag{$J_2$}\label{eq:J2}
\end{equation}
is the $p=\tfrac12$ specialization of~\eqref{eq:star}. Setting
$f(x) := G(x) - G(0)$, \eqref{eq:J2} is equivalent on a translate of~$I$
to Cauchy's additive equation $f(x+y) = f(x)+f(y)$ (cf.~\cite{Aczel1966}
or~\cite{Kuczma2009}). Cauchy additivity inherits a full
pathological-solution apparatus: without measurability, monotonicity,
boundedness on a set of positive measure, or another regularity hypothesis,
the equation admits non-affine solutions constructed via a \textbf{Hamel
basis} --- i.e., a $\Q$-vector-space basis of~$\R$, whose existence requires
the axiom of choice --- see~\cite{Hamel1905}. The classical theorems of
Cauchy~\cite{Cauchy1821}, Darboux~\cite{Darboux1875},
Hamel~\cite{Hamel1905}, Ostrowski~\cite{Ostrowski1929},
Sierpinski~\cite{Sierpinski1920}, and Steinhaus~\cite{Steinhaus1920}
together delineate which regularity hypotheses are enough to recover
affineness --- \emph{some} hypothesis is genuinely required, otherwise
affineness fails.

The \textbf{rational-coefficient} Jensen equation,
\begin{equation}
p\,G(u_1)+(1-p)\,G(u_2)=G(p\,u_1+(1-p)\,u_2)
\qquad (u_1,u_2\in I,\ p\in[0,1]\cap\Q),
\tag{$J_{\Q}$}\label{eq:JQ}
\end{equation}
sharpens \eqref{eq:J2} but does not escape its pathology: Cauchy additivity
yields $\Q$-homogeneity $f(qx) = qf(x)$ in the standard way ($G(nx)=nG(x)$
by induction, $G(x/n) = G(x)/n$ by substitution, $G(qx)=qG(x)$ for
$q\in\Q$ by combination), and $\Q$-homogeneity returns~\eqref{eq:JQ} for
$f$ and hence for~$G$. So \eqref{eq:JQ}, \eqref{eq:J2}, and Cauchy's
equation share the same solution class up to constants, and all three
inherit the Hamel-basis pathology.

The \textbf{continuous-coefficient} form~\eqref{eq:star}, by contrast, is
fundamentally different. It is the subject of this note.

\paragraph{Result.}
Suppose $G\colon[0,M]\to\R$ satisfies~\eqref{eq:star} for all
$u_1,u_2\in[0,M]$ and all $p\in[0,1]$. Then $G$ is affine on~$[0,M]$ ---
i.e., of the form $G(v)=av+b$ with $a,b\in\R$ --- explicitly,
\[
G(v) = a\,v + b, \qquad a = \frac{G(M)-G(0)}{M},\quad b=G(0).
\]
The proof is a single substitution: set $u_1=M$, $u_2=0$, $p=v/M$
in~\eqref{eq:star}. No regularity hypothesis on $G$ is used.

This is folklore --- at minimum~\cite{Aczel1966} and~\cite{Kuczma2009}
articulate the underlying observation, and the result is standard in the
functional-equations community. Nonetheless, an enduring confusion persists
in applied work that encounters~\eqref{eq:star} outside the
functional-equations literature: authors reach reflexively for the
Cauchy-equation regularity machinery --- continuity, boundedness,
measurability --- none of which the proof of affineness actually consumes.
The present note is a teaching-and-citation note: we exhibit the proof and
the dictionary of which hypotheses are vestigial for~\eqref{eq:star}, and
provide a one-line citation point for authors who encounter the equation in
their own work and wish to retire the Hamel concern without re-deriving it.

The structure is as follows. Section~\ref{sec:result} gives the theorem and
its proof. Section~\ref{sec:dictionary} is the dictionary of three regularity
hypotheses that are \emph{not} needed, each with the corresponding pathology
for~\eqref{eq:J2} that fails to arise for~\eqref{eq:star}.
Section~\ref{sec:variants} collects variants --- the weakest hypothesis the
proof actually consumes; the higher-dimensional remark; the contrast with
\eqref{eq:JQ}. Section~\ref{sec:recurrence} documents one recurrence of the
trap in the surrogate-calibration literature, together with the structural
reason the trap is predictable to recur in any derivation that pushes
Jensen's inequality to saturation.

% =============================================================================
\section{The result}
\label{sec:result}

\begin{theorem}
\label{thm:main}
Let $M>0$ and let $G\colon[0,M]\to\R$ satisfy~\eqref{eq:star} for all
$u_1,u_2\in[0,M]$ and all $p\in[0,1]$. Then $G$ is affine on~$[0,M]$:
\[
G(v) \;=\; \frac{G(M)-G(0)}{M}\,v \;+\; G(0)
\qquad \text{for all } v\in[0,M].
\]
No measurability, boundedness, or continuity hypothesis on $G$ is required.
\end{theorem}

\begin{proof}
Fix $v\in[0,M]$. Setting $u_1=M$, $u_2=0$, $p=v/M\in[0,1]$
in~\eqref{eq:star},
\[
\frac{v}{M}\,G(M) + \left(1-\frac{v}{M}\right)G(0)
\;=\; G\!\left(\frac{v}{M}\cdot M + \left(1-\frac{v}{M}\right)\cdot 0\right)
\;=\; G(v).
\]
Rearranging, $G(v) = G(0) + \dfrac{G(M)-G(0)}{M}\,v$.
\end{proof}

\begin{corollary}
\label{cor:regularity}
Under the hypotheses of Theorem~\ref{thm:main}, $G$ is in particular
continuous, monotone, locally Lipschitz, absolutely continuous, and
measurable on~$[0,M]$.
\end{corollary}

\begin{proof}
Affine functions are all of the above.
\end{proof}

The order of inference matters and is worth stating explicitly: the
\emph{conclusion} of Theorem~\ref{thm:main} is that $G$ has every regularity
property one might have wished to \emph{assume}. Each of those properties is
therefore vestigial as a hypothesis --- see Section~\ref{sec:dictionary}.

% =============================================================================
\section{The dictionary}
\label{sec:dictionary}

Set side-by-side the discrete-coefficient equation~\eqref{eq:J2} and the
continuous-coefficient equation~\eqref{eq:star}. The first column of
Table~\ref{tab:dictionary} names a regularity hypothesis that has been
classically invoked for~\eqref{eq:J2} and gives a one-sentence definition;
the second column records whether the hypothesis is required to force
affineness for~\eqref{eq:J2}, with the standard attribution; the third
column records what becomes of the hypothesis under~\eqref{eq:star}.

\begin{table}[ht]
\renewcommand{\arraystretch}{1.4}
\begin{tabularx}{\textwidth}{@{} X | X | X @{}}
\toprule
\textbf{Hypothesis on $G$ (definition)} & \textbf{Required for
\eqref{eq:J2}$\Rightarrow$ affine?} & \textbf{Required for
\eqref{eq:star}$\Rightarrow$ affine?}\\
\midrule
\textbf{Continuity on $I$} ($G$ is continuous at every $v\in I$).
& Yes, suffices (Cauchy~\cite{Cauchy1821}).
& \textbf{No} --- affineness is the \emph{conclusion}, not the hypothesis.\\
\midrule
\textbf{Measurability on $I$} ($G$ is Lebesgue- or Borel-measurable).
& Yes, suffices (Sierpinski~\cite{Sierpinski1920}).
& \textbf{No} --- same.\\
\midrule
\textbf{Monotonicity on $I$} ($G$ is non-decreasing or non-increasing).
& Yes, suffices (Darboux~\cite{Darboux1875}).
& \textbf{No} --- same.\\
\midrule
\textbf{Boundedness on a set of positive measure} (there exists a
Lebesgue-measurable $E\subseteq I$ with $|E|>0$ on which $G$ is bounded).
& Yes, suffices (Steinhaus~\cite{Steinhaus1920};
cf.~Sierpinski~\cite{Sierpinski1920}).
& \textbf{No} --- same.\\
\midrule
\textbf{Boundedness on $I$} ($G$ is bounded on the whole interval~$I$).
& Yes, suffices (special case of Steinhaus~\cite{Steinhaus1920};
Ostrowski~\cite{Ostrowski1929}).
& \textbf{No} --- same.\\
\midrule
\textbf{None} (no regularity hypothesis at all).
& \textbf{Insufficient} --- Hamel-basis pathology (Hamel~\cite{Hamel1905}).
& \textbf{Suffices} --- Theorem~\ref{thm:main}.\\
\bottomrule
\end{tabularx}
\caption{Regularity hypotheses on~$G$. Column~2 records which hypotheses
are classically required to force affineness for the discrete-coefficient
equation~\eqref{eq:J2}. Column~3 records what they become under the
continuous-coefficient equation~\eqref{eq:star}.}
\label{tab:dictionary}
\end{table}

The point of the table is the bottom row. \textbf{No regularity hypothesis}
is enough for the discrete-coefficient equation: a non-affine,
non-measurable, unbounded, additive function on~$\R$ --- the classical
Hamel-basis pathology of~\cite{Hamel1905} --- supplies a non-affine
solution of~\eqref{eq:J2}. By contrast, \textbf{no regularity hypothesis}
is \emph{needed} for the continuous-coefficient equation:
Theorem~\ref{thm:main}'s chord substitution closes the proof at the level
of the bare functional equation.

The mechanism is straightforward. A Hamel-basis pathological solution $G$
of Cauchy's equation on~$\R$ is $\Q$-linear: $G(0)=0$ and $G(qx) = qG(x)$
for every $q\in\Q$ and $x\in\R$. It is, by construction, \emph{not}
$\R$-linear --- $G(rx)\neq rG(x)$ for some irrational~$r$ and some
$x\in\R$.

The chord identity at the configuration $u_1=M$, $u_2=0$ reads
\begin{equation}
p\,G(M) + (1-p)\,G(0) \;=\; G(p\,M).
\tag{$\star_0$}\label{eq:star0}
\end{equation}
For $G$ a Hamel-pathological additive function ($G(0)=0$),
\eqref{eq:star0} is the assertion $pG(M) = G(pM)$ --- which holds at
$p\in\Q$ (by $\Q$-linearity) but fails at the irrational~$p$ where
$\R$-linearity is broken. So $G$ satisfies \eqref{eq:JQ} on $[0,M]$ but
violates~\eqref{eq:star} at \emph{every} irrational $p\in(0,1)$ where the
$\R$-linearity defect appears. The continuous-coefficient
equation~\eqref{eq:star} rules the pathology out at exactly the points
$p\in[0,1]\setminus\Q$ where the discrete- and rational-coefficient versions
are silent. \textbf{The Hamel pathology lives at irrational~$p$ --- exactly
where~\eqref{eq:star}'s continuous coefficient closes the door.}

% =============================================================================
\section{Variants and limits}
\label{sec:variants}

\subsection{The weakest hypothesis the proof consumes}
\label{ssec:weak}

The proof of Theorem~\ref{thm:main} uses~\eqref{eq:star} at exactly one
configuration: $u_1=M$, $u_2=0$, with $p\in[0,1]$ ranging freely. The full
force of~\eqref{eq:star} --- for \emph{every} pair $(u_1,u_2)$ --- is not
consumed. We therefore record the strict minimum.

\begin{theorem}[Theorem~\ref{thm:main}$'$]
\label{thm:main-prime}
Let $M>0$ and let $G\colon[0,M]\to\R$ satisfy
\[
p\,G(M) + (1-p)\,G(0) \;=\; G(p\,M)
\qquad \text{for all } p\in[0,1].
\]
Then $G(v) = G(0) + \bigl(G(M)-G(0)\bigr)\,v/M$ on $[0,M]$.
\end{theorem}

The proof is verbatim that of Theorem~\ref{thm:main}: set $p=v/M$.
Theorem~\ref{thm:main-prime} is in principle weaker than
Theorem~\ref{thm:main} --- the hypothesis~\eqref{eq:star0} does not a priori
imply the full~\eqref{eq:star}; under the conclusion (affineness) both hold.
In practice an author who \emph{derives}~\eqref{eq:star} from a richer setup
(e.g., from a Jensen-equality identity in a Bayes-risk computation, see
Section~\ref{sec:recurrence}) typically has~\eqref{eq:star} for all
$(u_1,u_2,p)$. The narrower Theorem~\ref{thm:main-prime} is the right
reference for an author who wants to know how much of the equation the proof
actually consumes.

\subsection{Convex domains in higher dimensions}
\label{ssec:higher}

The chord substitution extends verbatim to any convex subset~$C$ of a real
vector space~$V$: applying Theorem~\ref{thm:main} along every chord shows
that $G$ restricted to each segment of $C$ is affine in the segment
parameter, and a standard inductive argument
(cf.~\cite{AczelDhombres1989}) lifts chord-affineness to
convex-combination linearity on~$C$, hence to an affine form $G(x) =
a(x)+b$ for a linear~$a$ on $\mathrm{span}(C-C)$ and a constant $b\in\R$.
The structurally essential 1D step is Theorem~\ref{thm:main}.

\subsection{The rational-coefficient version $(J_{\Q})$ does retain the
pathology}
\label{ssec:rational-pathology}

The contrast across $\Q$ vs $\R$ is worth marking explicitly.

\begin{proposition}[folklore]
\label{prop:JQ-pathology}
There exists $G\colon[0,1]\to\R$ satisfying~\eqref{eq:JQ}
(and hence~\eqref{eq:J2}) on~$[0,1]$ that is not affine.
\end{proposition}

\begin{proof}[Construction]
Choose a Hamel basis~$H$ of~$\R$ over~$\Q$ containing~$1$. Define a
$\Q$-linear $\ell\colon\R\to\R$ by $\ell(1)=0$ and $\ell(h)=1$ for some
basis element $h\in H\setminus\{1\}$, and arbitrarily on the remaining
basis elements. By $\Q$-linearity, $\ell(x+y)=\ell(x)+\ell(y)$ for all
$x,y\in\R$ (so $\ell$ satisfies Cauchy's equation) and $\ell(qx)=q\ell(x)$
for all $q\in\Q$ and $x\in\R$ ($\Q$-homogeneity). But $\ell$ is \emph{not}
$\R$-linear: $\ell(1)=0$ while $\ell(h)=1$, so the values of $\ell$ on the
$\Q$-span of $\{1\}$ (namely $\Q$, on which $\ell\equiv 0$) and on the
$\Q$-span of $\{h\}$ (where $\ell$ is nonzero) are incompatible with any
$\R$-linear functional. Let $G := \ell|_{[0,1]}$. Then $G$
satisfies~\eqref{eq:JQ} on~$[0,1]$, since $\Q$-rational convex combinations
of points of $[0,1]$ stay in $[0,1]$ and $\ell$ respects $\Q$-linear
combinations on all of~$\R$.

Explicit witness that $G$ is not affine: choose $h\in H\setminus\{1\}$ with
$h\in(0,1)$ (such a basis element exists, e.g., $1/\pi$ extended to a
basis). Then $G(h)=\ell(h)=1\neq 0$, while $G(q)=\ell(q)=q\ell(1)=0$ for
every rational $q\in[0,1]$. Any affine map $A\colon[0,1]\to\R$ agreeing
with $G$ on $\Q\cap[0,1]$ would satisfy $A(0)=0$ and $A(q)=0$ for all
rational $q\in[0,1]$, forcing $A\equiv 0$ --- contradicting $G(h)=1$. So
$G$ is not affine.
\end{proof}

Proposition~\ref{prop:JQ-pathology} sharpens Theorem~\ref{thm:main}'s ``no
regularity hypothesis needed'' by exhibiting that the strengthening from
rational to continuous coefficients in~\eqref{eq:star} is doing real work.
Without it (i.e., for~\eqref{eq:JQ}), regularity hypotheses are genuinely
required; with it, they are vestigial.

% =============================================================================
\section{Where the trap recurs}
\label{sec:recurrence}

The chord substitution is --- as we have noted --- folklore. The regularity
hypotheses it makes vestigial nonetheless reappear in applied work that
derives~\eqref{eq:star} outside the functional-equations community. We
document \emph{one} such recurrence below --- in our own work on the
achievable error floor of partition-based classifiers --- and articulate
the \emph{structural source} under which the recurrence is predictable in
any sufficiently rich calibration-theory derivation. We invite extensions
of the catalog from readers who encounter~\eqref{eq:star} in their own work.

\subsection{A recurrence: surrogate calibration on the resolution axis}
\label{ssec:recurrence}

In recent work~\cite{El1,El2}, the equation~\eqref{eq:star} arises with~$G$
the function expressing the \emph{partition Bayes risk} of a measurable
classifier in terms of a concave score functional aggregated over the
partition's cells: $u_1,u_2$ are the per-cell score values and $p$ is the
cell mass, which ranges freely over $[0,1]$ on an atomless underlying
probability space, so that~\eqref{eq:star} holds in full. In the Lean~4
formalization of~\cite{El2}, the corresponding lemma was first declared
with a boundedness hypothesis in its signature, in deference to the
Cauchy/Hamel literature; the proof body then exhibited that the hypothesis
was unused, by exactly the chord substitution of Theorem~\ref{thm:main}.
The earlier drafts of the main text correspondingly invoked an apologetic
``\,$G$ is bounded so the bounded-Jensen-implies-affine theorem
applies\,'' parenthetical --- a parenthetical the present note's
Theorem~\ref{thm:main} retires.

\subsection{Structural source --- why the trap is predictable to recur}
\label{ssec:structural}

The recurrence is predictable. Whenever a calibration-theory argument
arrives at an identity of the form
\[
\E[g(\xi)] \;=\; g\bigl(\E[\xi]\bigr)
\qquad (\xi \text{ a random variable on } I,\ g\colon I\to\R),
\]
\emph{for a class of random variables $\xi$ wide enough that the marginal
$\E[\xi]$ can be any point of $I$ and the support of $\xi$ can be any
two-point subset of $I$ with any pair of masses $(p,1-p)$ with $p\in[0,1]$},
the identity is~\eqref{eq:star} with $g=G$, $\xi$ supported on
$\{u_1,u_2\}$. Such an identity arises whenever Jensen's inequality is
pushed to saturation. The standard surrogate-calibration literature ---
Bartlett--Jordan--McAuliffe~\cite{BJM2006},
Tewari--Bartlett~\cite{TewariBartlett2007},
Steinwart~\cite{Steinwart2007}, Reid--Williamson~\cite{ReidWilliamson2010,
ReidWilliamson2011} --- derives its calibration results via convex
analysis (biconjugates, Fenchel duality, supporting hyperplanes), which
keeps Jensen as an \emph{inequality} with explicitly controlled slack and
consequently never produces~\eqref{eq:star} as an equation; the
resolution-axis derivation in~\cite{El2} pushes Jensen to equality directly
and thereby produces~\eqref{eq:star}. The chord substitution closes the
resulting equation in one line. Other derivation styles that saturate
Jensen's inequality across a wide class of two-point distributions should
be expected to produce~\eqref{eq:star} similarly --- at which point
Theorem~\ref{thm:main} retires the regularity-hypothesis question on the
spot.

\subsection{Adjacent settings}
\label{ssec:adjacent}

Two adjacent settings where structurally similar equations arise and the
chord substitution is the natural closure, although the \emph{trap}
(invoking Cauchy/Hamel regularity unnecessarily) is not, to our knowledge,
recurrent in published work. First, \textbf{expected-utility representation
theorems} in the von~Neumann--Morgenstern tradition derive the
linearity-in-probability of a utility functional
$U(p\,L_1+(1-p)\,L_2)=p\,U(L_1)+(1-p)\,U(L_2)$ on lotteries; the
Herstein--Milnor axiomatization closes this via the \emph{Archimedean
axiom} rather than via Cauchy/Hamel, but the chord substitution is the
natural alternative closure. Second, \textbf{Shannon entropy's axiomatic
characterization} (Khinchin--Faddeev recursivity axiom; cf.~\cite{Aczel1966}
and~\cite{AczelDhombres1989}) treats richer functional equations
than~\eqref{eq:star}, but several intermediate steps reduce to
continuous-coefficient Jensen-type identities for which the chord
substitution is part of the standard toolkit. We flag these settings as
adjacent rather than as documented instances of the trap.

\subsection{An invitation to extend}
\label{ssec:invitation}

We invite readers who have encountered~\eqref{eq:star} in their own work
and have invoked a vestigial regularity hypothesis to extend the catalog
with a one-line citation back to the relevant proof step in their paper.

% =============================================================================
\section*{Declaration of interests}

The author does not work for, advise, own shares in, or receive funds from
any organisation that could benefit from this article, and has declared no
affiliation other than his/her research organisation.

% =============================================================================
\section*{Acknowledgements}

The chord-substitution observation surfaced in the Lean~4 formalization
track of~\cite{El2}, during a phase of the proof development whose
discipline is captured in three open project skills documented at the
source code repository associated with~\cite{El2}. The author thanks the
internal adversarial-review process described in the same source for
surfacing the over-engineered boundedness hypothesis whose retirement
motivated the present note.

% =============================================================================
% Bibliography
% =============================================================================

\bibliographystyle{plain} % FALLBACK: replace with `crmath` style when
                          % the official Centre Mersenne class is installed
\bibliography{refs}

\end{document}
